\chapter{Conclusion}
\label{ch:conclusion}

This thesis set out to make visual drone detection deployable: high recall and a low false-alarm rate at the same time, on hardware a site can afford. The position it defends is architectural: the base detectors are tuned purely for recall, and every discrimination they cannot learn without losing recall (confuser rejection, modality arbitration, false-positive filtering) is relocated into cheap, learned, downstream stages. The evaluation holds one protocol throughout (Section~\ref{sec:eval_protocol}); every number below carries its configuration and its confidence interval in the tables it cites.

\section{Production Stack}
\label{sec:production_stack}

The production stack, as validated in this thesis, is as follows. The RGB detector is the SelCom-plus-confuser fine-tune, the only confuser-injection recipe that passed every drone-recall regression gate. The thermal detector is the six-revision model of Section~\ref{sec:ir_evolution}. The trust router uses eight free features, selected by the leakage statistic, and always routes (to RGB, IR, or both), leaving false-positive rejection to the filter; the reject-capable variant is kept as an ablation and is the better choice on confuser-rich paired streams. Each modality carries its own confuser filter, run per frame: an RGB filter at threshold $0.25$ and a thermal-native IR filter at $0.05$. An $N$-of-$M$ temporal smoother sits above them. Inference runs at \texttt{imgsz=1280} on Svanstr\"om-like and CCTV content (960 is the measured optimum for the deployment camera) and at 640 elsewhere, and trust-aware scoring is the reporting rule. The exact weights for each stage are listed in Appendix~\ref{app:models}. One validated operating mode extends the stack: on recall-starved out-of-distribution RGB surfaces, the RGB confidence floor can drop to $0.05$--$0.10$ and let the filter own precision, since that is where false-positive control now lives ($+10$~pp F1 on SelCom at unchanged confuser safety; Section~\ref{sec:lowconf_mode}).

Two carve-outs accompany it, each diagnosed rather than hidden. The earlier coverage gap on the in-domain RGB test split, where the confuser filter once cost $11$~pp F1, is now closed: the bird-split rebuild lifts that split's F1 from $0.792$ to $0.916$ (recall $0.691 \to 0.887$), so the patch filter is no longer needed there as a fallback (Section~\ref{sec:verifier_results}). What remains is a deliberate trade. First, the shipped router gives up some confuser suppression for recall: RGB-confuser frame fire rises to about $1.4\%$, and real-video-confuser fire roughly triples, relative to the reject-capable router, which is kept as the recommended choice where confuser-rich paired video dominates (Section~\ref{sec:classifier_results}). Second, the thermal-native IR filter is validated per frame in the evaluation engine, but its operator-GUI wiring is an open engineering item at the time of writing, as is Jetson-class edge latency.

\section{Future Work}

Four directions. A \emph{track-based classifier} consuming the temporal smoother's per-track state (velocity, persistence, trajectory curvature) as a fourth evidence channel, the most promising single direction, and the one aimed at the airplane class that appearance-based stages handle least well. \emph{Further hardening the airplane class}: the thermal-native filter already cut the airplane-dominated thermal-confuser fire $29.4\% \to 2.8\%$ (cache) / $94\%$ (held-out), and additional OOD airplane crops as confuser-class training data would press the residual lower. \emph{Completing the engineering}: GUI wiring of the aligned filter and TensorRT/ONNX edge benchmarking. And \emph{extending the protocol}: the conservative open-world router remains selectable where missed drones are operationally cheaper than false alarms, and the dual-classifier extension (Section~\ref{sec:classifier_results}) is validated if a deployment needs the missing-modality regimes as first-class citizens.
